\chapter{Supplementary Figures and Tables}
\label{append:comparisons}
\vspace*{2em}

This appendix holds reference material summarised in the body: the
deployment models considered in \autoref{design:sec:deployment}, the cloud
providers considered in \autoref{impl:subsec:provider}, the inference
configurations described in \autoref{impl:subsec:configs}, the application
components described in \autoref{impl:sec:stack}, and the instructor-hours
projection of \autoref{eval:subsec:hours}.

\begin{table}[htb]
  \centering
  \caption{Comparison of deployment models for the virtual teaching assistant.}
  \label{design:tab:deployment}
  \footnotesize
  \renewcommand{\arraystretch}{1.3}
  \begin{tabular}{L{0.13\textwidth} L{0.18\textwidth} L{0.18\textwidth} L{0.18\textwidth} L{0.185\textwidth}}
    \hline
    \textbf{Dimension} & \textbf{Cloud} & \textbf{On-Premises} & \textbf{Hybrid} & \textbf{Relevance to This System} \\
    \hline
    Scalability &
      Auto-scaling, elastic resources &
      Fixed hardware capacity &
      Burst to cloud for peak load &
      Needed when many students interact simultaneously \\
    Initial cost (CapEx) &
      Low upfront cost &
      High hardware investment &
      Moderate; reuses existing hardware &
      Important for student projects with limited budget \\
    Operational cost (OpEx) &
      Pay-as-you-go pricing &
      Maintenance, electricity, upgrades &
      Both cost models apply &
      Cloud easier to forecast for experimentation \\
    AI/ML infrastructure &
      Managed GPUs, LLM APIs, training services &
      GPUs and frameworks configured manually &
      Local inference; managed services for heavy workloads &
      Faster iteration for adaptive learning models \\
    Maintenance and updates &
      Managed by provider &
      Self-managed patches and upgrades &
      Split responsibility &
      Frees effort for algorithm design over operations \\
    Security controls &
      Built-in IAM, encryption, monitoring &
      Full control, higher responsibility &
      Must secure both environments and the link &
      Student learning data requires strong protection \\
    Data privacy and control &
      Shared responsibility model &
      Full ownership of data &
      Sensitive data on-premises; compute in cloud &
      On-premises stronger for sensitive datasets \\
    Reliability and availability &
      Multi-region redundancy &
      Depends on local setup &
      Cloud fallback for local outages &
      Important for continuous student access \\
    Latency &
      Depends on region selection &
      Very low on local network &
      Low for local paths; variable for cloud paths &
      Useful for real-time AI responses \\
    Development speed &
      Fast prototyping with managed services &
      Longer setup time &
      Added integration complexity &
      Accelerates the project timeline \\
    Vendor lock-in risk &
      Higher (managed services) &
      Lower &
      Lower; workloads remain portable &
      Key architectural consideration \\
    \hline
  \end{tabular}
\end{table}

\begin{table}[htb]
  \centering
  \caption{Comparison of major cloud providers for the virtual teaching
           assistant.}
  \label{impl:tab:providers}
  \footnotesize
  \renewcommand{\arraystretch}{1.3}
  \begin{tabular}{L{0.13\textwidth} L{0.18\textwidth} L{0.18\textwidth} L{0.18\textwidth} L{0.185\textwidth}}
    \hline
    \textbf{Dimension} & \textbf{AWS} & \textbf{Google Cloud (GCP)} & \textbf{Microsoft Azure} & \textbf{Relevance to This System} \\
    \hline
    Managed AI services &
      Bedrock, SageMaker &
      Vertex AI (strong ML workflows) &
      Azure OpenAI Service &
      Needed for question generation and adaptive learning \\
    LLM integration &
      Broad model ecosystem &
      Strong data--ML integration &
      Deep OpenAI ecosystem &
      Influences development complexity \\
    Data analytics stack &
      Redshift, Kinesis &
      BigQuery, Pub/Sub (very strong) &
      Synapse, Event Hubs &
      Useful for tracking student metrics and analytics \\
    Serverless compute &
      Lambda &
      Cloud Functions &
      Azure Functions &
      Suits event-driven microservices \\
    Kubernetes support &
      EKS &
      GKE (very mature) &
      AKS &
      Relevant for a microservices architecture \\
    Developer experience &
      Very large ecosystem &
      Clean ML workflows &
      Strong enterprise integration &
      Impacts development speed \\
    Identity and access management &
      Powerful but complex &
      Simpler structure &
      Tight Microsoft-ecosystem integration &
      Security for student accounts \\
    Pricing model &
      Flexible but complex &
      Often cost-efficient for analytics &
      Enterprise-oriented pricing &
      Important for long-term deployment \\
    Learning curve &
      Medium--high &
      Lower for ML-focused projects &
      Medium &
      Relevant for a solo developer \\
    Best strength &
      Largest ecosystem and flexibility &
      Data and AI workflow excellence &
      OpenAI integration and enterprise tools &
      Helps justify the final platform choice \\
    \hline
  \end{tabular}
\end{table}

\begin{figure}[htb]
  \centering
  % Application components, in layers. Simplified from
% cp5105-master-architecture.drawio, View 2 (per-feature wiring omitted),
% with the Groq provider in place of Gemini.
\begin{tikzpicture}[dg]
  \def\cA{1.55} \def\cB{4.4} \def\cC{7.25} \def\cD{10.1} \def\cE{12.95}

  % browser
  \foreach \x/\n/\t/\f in {\cA/u1/Sign-in/auth.js, \cB/u2/Courses (LMS)/lms.js,
      \cC/u3/Chat/chat.js, \cD/u4/Quiz/quiz.js, \cE/u5/Dashboards/stats.js}
    \node[box=25mm, blue] (\n) at (\x,0) {\textbf{\t}\\ \texttt{\f}};
  \begin{scope}[on background layer]
    \node[grp, fit=(u1)(u5), inner ysep=3mm, yshift=1.5mm] (b1) {};
  \end{scope}
  \node[grplabel] at (b1.north west) {Browser: plain JavaScript modules over one shared state};

  % HTTP boundary
  \node[box=53mm] (ngx) at ($(\cA,-1.95)!0.5!(\cB,-1.95)$)
    {\textbf{Nginx}: serves the app, proxies \texttt{/api}};
  \node[box=81mm] (gate) at ($(\cC,-1.95)!0.5!(\cE,-1.95)$)
    {\textbf{Auth gate}: session cookie and role on every \texttt{/api} route};
  \begin{scope}[on background layer]
    \node[grp, fit=(ngx)(gate), inner ysep=3mm, yshift=1.5mm] (b2) {};
  \end{scope}
  \node[grplabel] at (b2.north west) {HTTP boundary};

  % request handlers
  \foreach \x/\n/\t/\f in {\cA/h1/Auth/auth.js, \cB/h2/LMS connectors/{canvas, classroom},
      \cC/h3/Chat agent/chat.js, \cD/h4/Adaptive quiz/quiz.js, \cE/h5/Analytics/stats.js}
    \node[box=25mm, violet] (\n) at (\x,-3.9) {\textbf{\t}\\ \texttt{\f}};
  \begin{scope}[on background layer]
    \node[grp, fit=(h1)(h5), inner ysep=3mm, yshift=1.5mm] (b3) {};
  \end{scope}
  \node[grplabel] at (b3.north west) {Request handlers: one module per feature};

  % shared services (two rows)
  \foreach \x/\n/\t/\f in {2.0/s1/Retrieval/rag.js, 5.6/s2/Embeddings/embeddings.js,
      9.2/s3/Agent tools/agentTools.js, 12.8/s4/MCP client/mcpClient.js}
    \node[box=31mm, green] (\n) at (\x,-5.85) {\textbf{\t}\\ \texttt{\f}};
  \foreach \x/\n/\t/\f in {2.0/s5/LLM router/{llm/index.js}, 5.6/s6/SQL access/db\_index.js,
      9.2/s7/Fuzzy feedback/feedbackEngine.py, 12.8/s8/Web search/search.js}
    \node[box=31mm, green] (\n) at (\x,-6.95) {\textbf{\t}\\ \texttt{\f}};
  \begin{scope}[on background layer]
    \node[grp, fit=(s1)(s8), inner ysep=3mm, yshift=1.5mm] (b4) {};
  \end{scope}
  \node[grplabel] at (b4.north west) {Shared services: each loads on its own; a failure disables only its feature};

  % data and external systems
  \foreach \x/\n/\t/\f in {1.25/x1/Postgres/+ pgvector, 3.65/x2/Ollama/local, 6.05/x3/Groq/hosted,
      8.45/x4/{Canvas,}/Classroom, 10.85/x5/MCP/servers, 13.25/x6/Web search/API}
    \node[box=19mm, amber] (\n) at (\x,-8.8) {\textbf{\t}\\ \f};
  \begin{scope}[on background layer]
    \node[grp, fit=(x1)(x6), inner ysep=3mm, yshift=1.5mm] (b5) {};
  \end{scope}
  \node[grplabel] at (b5.north west) {Data and external systems};

  % layer-to-layer flow
  \foreach \a/\b/\l in {b1/b2/\texttt{fetch('/api/...')}, b2/b3/dispatch by route,
      b3/b4/calls, b4/b5/{SQL, HTTP}}
    \draw[flow] (\a.south -| \cC,0) ++(1.6,0) -- node[lbl, right=1pt] {\l} ($(\b.north -| \cC,0)+(1.6,0)$);
\end{tikzpicture}

  \caption{Application components, by layer. Each feature has one browser
           module and one request handler; the handlers share the services
           in the fourth layer.}
  \label{impl:fig:components}
\end{figure}

\begin{table}[htb]
  \centering
  \caption{Inference configurations used during the project.}
  \label{impl:tab:configs}
  \footnotesize
  \renewcommand{\arraystretch}{1.3}
  \begin{tabular}{L{0.14\textwidth} L{0.37\textwidth} L{0.37\textwidth}}
    \hline
    \textbf{Period} & \textbf{Configuration} & \textbf{Results measured} \\
    \hline
    Development &
      Fully local: Llama~3.2 (3B) via Ollama on an 8\,GB consumer laptop &
      Compression probe (\autoref{eval:sec:results:compression}) \\
    Aug 2026 &
      Fully local: Qwen3 14B (agent) and Llama~3.2 (quiz) via Ollama on a
      32\,GB M1 Max, reached from EC2 over a reverse SSH tunnel &
      Local rows of the provider benchmark
      (\autoref{eval:sec:results:providers}) \\
    From 27 Sep 2026 &
      Hybrid: gpt-oss-20b (agent) and gpt-oss-120b (quiz) on Groq; BGE-M3
      embeddings and Llama~3.2 fallback local &
      Hosted rows of the provider benchmark; token audit
      (\autoref{design:subsec:tokenaudit}); answer-quality evaluation
      (\autoref{eval:sec:results:quality}); cost projection
      (\autoref{eval:subsec:results:cost}) \\
    \hline
  \end{tabular}
\end{table}

\begin{table}[htb]
  \centering
  \caption{Illustrative projection of instructor-hours returned per course per
           week (\autoref{eval:subsec:hours}). Every scenario uses the measured
           65\% fully-correct rate; all other values are assumptions.}
  \label{append:tab:hours}
  \footnotesize
  \renewcommand{\arraystretch}{1.25}
  \begin{tabular}{L{0.42\textwidth} r r r}
    \hline
    \textbf{Parameter} & \textbf{Low} & \textbf{Central} & \textbf{High} \\
    \hline
    Enrolment, $N$                                         & 200  & 250  & 300  \\
    Questions per student per week, $q$\textsuperscript{a} & 0.14 & 0.36 & 0.59 \\
    Share of questions that are routine, $r$               & 60\% & 75\% & 90\% \\
    Share answered fully correctly, $d$ (measured)         & 65\% & 65\% & 65\% \\
    Share that would otherwise reach staff, $s$            & 20\% & 35\% & 50\% \\
    Look-up time per question (min)                        & 0.5  & 1.0  & 2.0  \\
    Reply length (words), composed at 19 wpm\textsuperscript{b} & 40 & 75 & 120 \\
    Handling time per question, $t$ (min)\textsuperscript{c} & 2.7 & 5.1 & 8.4 \\
    \hline
    Questions taken off staff per week ($N q r d s$)        & 2.2  & 15.4 & 51.4 \\
    \textbf{Instructor-hours per course per week}          & \textbf{0.1} & \textbf{1.3} & \textbf{7.2} \\
    Instructor-hours per 13-week semester                  & 1    & 17   & 94   \\
    \hline
\end{tabular}

  \vspace{4pt}
  \parbox{0.86\textwidth}{\scriptsize
    \textsuperscript{a}KAIST deployment~\cite{acl25:Kweon}: 916 conversations
    (low) or 3\,869 exchanges (high) from 472 students over 14 weeks; central is
    their midpoint.
    \textsuperscript{b}Composition typing rate reported by Karat et
    al.~\cite{chi99:Karat}.
    \textsuperscript{c}Reading a 25-word question at 238 words per
    minute~\cite{jml19:Brysbaert}, plus look-up, plus composing the reply.}
\end{table}
